% Options for packages loaded elsewhere
% Options for packages loaded elsewhere
\PassOptionsToPackage{unicode}{hyperref}
\PassOptionsToPackage{hyphens}{url}
\PassOptionsToPackage{space}{xeCJK}
%
\documentclass[
  letterpaper,
  oneside]{scrbook}
\usepackage{xcolor}
\usepackage[margin=20mm]{geometry}
\usepackage{amsmath,amssymb}
\setcounter{secnumdepth}{-\maxdimen} % remove section numbering
\usepackage{iftex}
\ifPDFTeX
  \usepackage[T1]{fontenc}
  \usepackage[utf8]{inputenc}
  \usepackage{textcomp} % provide euro and other symbols
\else % if luatex or xetex
  \usepackage{unicode-math} % this also loads fontspec
  \defaultfontfeatures{Scale=MatchLowercase}
  \defaultfontfeatures[\rmfamily]{Ligatures=TeX,Scale=1}
\fi
\usepackage{lmodern}
\ifPDFTeX\else
  % xetex/luatex font selection
  \setmainfont[]{Times New Roman}
  \ifXeTeX
    \usepackage{xeCJK}
    \setCJKmainfont[]{meiryo.ttc}
  \fi
  \ifLuaTeX
    \usepackage[]{luatexja-fontspec}
    \setmainjfont[]{meiryo.ttc}
  \fi
\fi
% Use upquote if available, for straight quotes in verbatim environments
\IfFileExists{upquote.sty}{\usepackage{upquote}}{}
\IfFileExists{microtype.sty}{% use microtype if available
  \usepackage[]{microtype}
  \UseMicrotypeSet[protrusion]{basicmath} % disable protrusion for tt fonts
}{}
\makeatletter
\@ifundefined{KOMAClassName}{% if non-KOMA class
  \IfFileExists{parskip.sty}{%
    \usepackage{parskip}
  }{% else
    \setlength{\parindent}{0pt}
    \setlength{\parskip}{6pt plus 2pt minus 1pt}}
}{% if KOMA class
  \KOMAoptions{parskip=half}}
\makeatother
% Make \paragraph and \subparagraph free-standing
\makeatletter
\ifx\paragraph\undefined\else
  \let\oldparagraph\paragraph
  \renewcommand{\paragraph}{
    \@ifstar
      \xxxParagraphStar
      \xxxParagraphNoStar
  }
  \newcommand{\xxxParagraphStar}[1]{\oldparagraph*{#1}\mbox{}}
  \newcommand{\xxxParagraphNoStar}[1]{\oldparagraph{#1}\mbox{}}
\fi
\ifx\subparagraph\undefined\else
  \let\oldsubparagraph\subparagraph
  \renewcommand{\subparagraph}{
    \@ifstar
      \xxxSubParagraphStar
      \xxxSubParagraphNoStar
  }
  \newcommand{\xxxSubParagraphStar}[1]{\oldsubparagraph*{#1}\mbox{}}
  \newcommand{\xxxSubParagraphNoStar}[1]{\oldsubparagraph{#1}\mbox{}}
\fi
\makeatother


\usepackage{longtable,booktabs,array}
\newcounter{none} % for unnumbered tables
\usepackage{calc} % for calculating minipage widths
% Correct order of tables after \paragraph or \subparagraph
\usepackage{etoolbox}
\makeatletter
\patchcmd\longtable{\par}{\if@noskipsec\mbox{}\fi\par}{}{}
\makeatother
% Allow footnotes in longtable head/foot
\IfFileExists{footnotehyper.sty}{\usepackage{footnotehyper}}{\usepackage{footnote}}
\makesavenoteenv{longtable}
\usepackage{graphicx}
\makeatletter
\newsavebox\pandoc@box
\newcommand*\pandocbounded[1]{% scales image to fit in text height/width
  \sbox\pandoc@box{#1}%
  \Gscale@div\@tempa{\textheight}{\dimexpr\ht\pandoc@box+\dp\pandoc@box\relax}%
  \Gscale@div\@tempb{\linewidth}{\wd\pandoc@box}%
  \ifdim\@tempb\p@<\@tempa\p@\let\@tempa\@tempb\fi% select the smaller of both
  \ifdim\@tempa\p@<\p@\scalebox{\@tempa}{\usebox\pandoc@box}%
  \else\usebox{\pandoc@box}%
  \fi%
}
% Set default figure placement to htbp
\def\fps@figure{htbp}
\makeatother





\setlength{\emergencystretch}{3em} % prevent overfull lines

\providecommand{\tightlist}{%
  \setlength{\itemsep}{0pt}\setlength{\parskip}{0pt}}



 


\makeatletter
\@ifpackageloaded{bookmark}{}{\usepackage{bookmark}}
\makeatother
\makeatletter
\@ifpackageloaded{caption}{}{\usepackage{caption}}
\AtBeginDocument{%
\ifdefined\contentsname
  \renewcommand*\contentsname{Table of contents}
\else
  \newcommand\contentsname{Table of contents}
\fi
\ifdefined\listfigurename
  \renewcommand*\listfigurename{List of Figures}
\else
  \newcommand\listfigurename{List of Figures}
\fi
\ifdefined\listtablename
  \renewcommand*\listtablename{List of Tables}
\else
  \newcommand\listtablename{List of Tables}
\fi
\ifdefined\figurename
  \renewcommand*\figurename{Figure}
\else
  \newcommand\figurename{Figure}
\fi
\ifdefined\tablename
  \renewcommand*\tablename{Table}
\else
  \newcommand\tablename{Table}
\fi
}
\@ifpackageloaded{float}{}{\usepackage{float}}
\floatstyle{ruled}
\@ifundefined{c@chapter}{\newfloat{codelisting}{h}{lop}}{\newfloat{codelisting}{h}{lop}[chapter]}
\floatname{codelisting}{Listing}
\newcommand*\listoflistings{\listof{codelisting}{List of Listings}}
\makeatother
\makeatletter
\makeatother
\makeatletter
\@ifpackageloaded{caption}{}{\usepackage{caption}}
\@ifpackageloaded{subcaption}{}{\usepackage{subcaption}}
\makeatother
\usepackage{bookmark}
\IfFileExists{xurl.sty}{\usepackage{xurl}}{} % add URL line breaks if available
\urlstyle{same}
\hypersetup{
  pdftitle={Notes from a Moving Silence - 移りゆく静寂からの手記},
  pdfauthor={David Tord - デイヴィッド・トールド},
  hidelinks,
  pdfcreator={LaTeX via pandoc}}


\title{Notes from a Moving Silence - 移りゆく静寂からの手記}
\author{David Tord - デイヴィッド・トールド}
\date{2026-01-01}
\begin{document}
\frontmatter
\maketitle


\mainmatter
\bookmarksetup{startatroot}

\chapter{Preface}\label{preface}

\bookmarksetup{startatroot}

\chapter{序文}\label{ux5e8fux6587}

These notes would not exist without the people who made this journey
possible, who welcomed a stranger, shared their stories, and transformed
what began as a simple visit into something far more meaningful.
この記録は、この旅を可能にしてくれた人々がいなければ存在しなかっただろう。彼らはよそ者である私を迎え入れ、物語を分かち合い、単なる訪問から始まったものを、はるかに意味深いものへと変えてくれた。

To Morino-san, Miki-san, Yamato-kun and Sakura-chan thanks for your
kindness and patience as I fumbled through conversations and culture set
the foundation for everything that followed. Thank you for seeing past
the language barriers and for your generosity.
森野さん、ミキさん、ヤマトくん、そしてサクラちゃんへ。言葉や文化に手探り状態だった私に、親切に根気よく接してくれたことが、その後のすべての土台となりました。言葉の壁を越えて向き合ってくれたこと、そしてその寛大さに感謝します。

Kobayashi-san, you were the quiet teacher. Every day, a lesson in how to
live. How to work. How to move through the world with grace, curiosity,
and genuine care. You guided these notes simply by being yourself.
小林さん、あなたは静かなる師でした。毎日が、どう生きるか、どう働くかの学びでした。優雅さと好奇心、そして心からの思いやりを持って、この世界をどう進んでいくか。あなたは、ただあなた自身でいることによって、この記録を導いてくれました。

To the Warabiya Kids and Akiko-san, you reminded me what it means to be
fully present. Whether it was fishing in the rain, learning the stories
of your island, or simply existing in the rhythm of your home, you
taught me more than you know.
わらびやの子供たちとアキコさんへ。皆さんは、今この瞬間に完全に意識を向けることの意味を思い出させてくれました。雨の中での釣りも、島の物語を学ぶことも、あるいは単にあなたたちの家のリズムの中で過ごすことも、私がどれほど多くを学んだか、皆さんは気付いていないかもしれません。

Rai-san, your quiet strength left an impression that lingers still.
ライさん、あなたの静かなる強さは、今もなお私の中に深い印象を残しています。

To Yoko sensei and Michael, thank you for the conversations, the
laughter, and for helping me understand the deeper currents running
through this place.
ヨーコ先生とマイケルへ。たくさんの語らいと笑いを、そしてこの場所に流れる深い源流を理解する手助けをしてくれて、本当にありがとう。

To Mai-chan, your spirit and joy were a light during the quieter
moments.
マイちゃん、あなたの明るい精神と喜びは、静かな時間の中で私を照らす光でした。

And to everyone on the island, the fishermen at the port, the families
we crossed paths with, the elders whose presence carried the weight of
history, and the children whose curiosity mirrored my own. You are woven
into every page of this book. Your generosity, your resilience, and the
way you move through life with such graceful intentionality, that is
what these notes attempt to capture.
そして島のみんなへ。港の漁師たち、すれ違った家族、その存在が歴史の重みを物語っていたお年寄りたち、そして私と同じ好奇心を抱いていた子供たち。皆さんは、この本のすべてのページに織り込まれています。皆さんの寛大さ、たくましさ、そして優雅な意志を持って人生を歩むその姿こそが、私がこの記録に留めようとしたものです。

This collection is not a guide. It is a love letter to a place and the
people who live there, written in fragments, in the quiet spaces between
arrival and departure. It is my attempt to hold onto something I cannot
stay in, and to share it with others who might understand that
sometimes, the most profound experiences come from simply being present;
watching clouds move across the sea, listening to the rain, and allowing
ourselves to be changed by the people we meet along the way.
この作品集はガイドブックではありません。これは、ある場所とそこに住む人々へのラブレターであり、到着から出発までの静かな空間の中で、断片的に書かれたものです。これは、私が永遠に留まることのできない場所での記憶を抱きしめ、共有するための試みです。時には、ただそこにいること、つまり海を渡る雲を眺め、雨音に耳を澄ませ、道中で出会う人々によって自分が変わっていくのに身を任せることからこそ、最も深い経験が生まれるのだということを理解してくれる人たちに向けて。

\begin{center}\rule{0.5\linewidth}{0.5pt}\end{center}

\bookmarksetup{startatroot}

\chapter{Notes from a Moving Silence}\label{notes-from-a-moving-silence}

I spent most of my free time at Tokashiku Beach. A place that felt
almost unreal.

A place where birds crossed paths with whales. Where octopuses met the
crescent moon. Where the foam of the ocean blended into the clouds.

Those clouds I would watch for hours, lying on the white remains of
parrotfish, scattered across the sand.

We gave them names. Or at least, we tried to. It didn't matter much.

What struck me most was their speed. They moved across the East China
Sea faster than anything I had ever seen , even faster than those pesky
F-16 jets.

But lately, it was the opposite kind that held my attention.

Slow, heavy clouds. Thick, almost motionless. They could stay there for
hours without changing, without drifting.

They weren't threatening. They weren't even particularly beautiful.

They were just there.

Waiting, quietly, like notes from a moving silence..

\begin{center}\rule{0.5\linewidth}{0.5pt}\end{center}

渡嘉敷ビーチで、私は自由な時間のほとんどを過ごしていた。どこか現実離れした場所だった。

鳥とクジラがすれ違い、タコが三日月と出会い、海の泡が雲に溶けていくような場所。

砂浜に散らばるブダイの白い痕の上に寝転びながら、私は何時間も雲を眺めていた。

私たちは雲に名前をつけようとした。少なくとも、そうしようとはしていた。けれど、あまり重要なことではなかった。

いちばん印象に残ったのは、その速さだった。東シナ海を横切っていくその動きは、今まで見た何よりも速くて、あのやかましいF16戦闘機よりもずっと速く感じられた。

でも最近、私の目を引くのは、その反対のものだった。

ゆっくりとした、重たい雲。厚く、ほとんど動かない。何時間も、そのままの形で、漂うことすらなくそこに留まっている。

威圧的でもない。特別に美しいわけでもない。

ただ、そこにあるだけ。

静かに、ただ待っている。動いているはずの静けさの中に浮かぶ音符のように。

\bookmarksetup{startatroot}

\chapter{Part I - Fishing Through the
Storm}\label{part-i---fishing-through-the-storm}

I used to see the kids every day, running, biking, fishing rods over
their shoulders, rushing out after school, always heading somewhere
different. I kept wondering how news of the \emph{good} fishing spots
traveled so fast between them. Some kind of island-grade, low-tech,
child-powered communication network.

One afternoon, they gathered at the port. It happened to be my day off.
So I jumped on my tiny pink, aggressively kawaii folding bike, grabbed
my rod, and rode back to join them, pretending this was all very casual.

They were shy at first. Quick smiles, eyes darting back to their lines.
A few were fooling around, but most were intensely focused, serious in
that way only kids can be when something actually matters.

Then Yuma showed up. He's Rai-san's son. We'd already bonded over
Beyblades and Pokémon at Sunny Coral, which, as far as friendships go,
is a solid foundation. He walked straight up to me and explained, very
carefully, how pier fishing actually works. I'd mostly fished rivers
back in the Swiss mountains. This was a different world.

\begin{figure}[H]

{\centering \pandocbounded{\includegraphics[keepaspectratio]{.attachments/ea8ecd5cf58f5362337eb67873d16ff7fd0da431.jpeg}}

}

\caption{image}

\end{figure}%

He asked me my favorite Pokémon. I wanted to say \emph{Scyther}, but
didn't know the Japanese name, so I told a small, strategic lie.

``\emph{Lizardon ga suki desu.}''

Charizard.

He nodded vigorously. Approval granted. Then, realizing I could
understand \emph{some} Japanese, as kids and elderly people love to do,
he immediately switched to full-speed mode. A verbal flamethrower. I
nodded back, just as vigorously, hoping enthusiasm could replace
comprehension.

Rain started as a drizzle. No one moved. Then it turned into real rain.
Still, no one moved. Finally, it became proper Okinawan rain,
horizontal, loud, soaking, disrespectful.

\begin{figure}[H]

{\centering \pandocbounded{\includegraphics[keepaspectratio]{.attachments/3c488e839ef6116b2433b40925363895c6a304d7.jpeg}}

}

\caption{image}

\end{figure}%

Even with my Gore-Tex jacket, I was drenched, so I just gave up to the
rain. Some mums came to retrieve their kids. I waved goodbye to Yuma. A
small group remained, completely unbothered, rods still in the water.

That's when Akiko-san appeared, carrying umbrellas. Not to stop them,
she clearly knew better, but to make the situation survivable. She
smiled at me, tall and calm, the kind of presence that makes chaos feel
organized. In flawless English, she joked that I was a very motivated
fisherman.

We talked in the rain. I explained I was helping at the guesthouse. She
told me the kids weren't hers --- they lived at her place, Warabiya, a
kind of boarding home for children from all over Japan. Tokyo, Osaka.
Parents sending them here to grow up close to the sea and forest,
learning by doing. I thought back to the Protestant boarding school in
Fribourg I had narrowly escaped as a kid. Same concept. Very different
execution.

That's when the sign I'd seen near my house finally made sense
(わらびや).

After half an hour of rain-soaked conversation, we headed home. It was
my first day with the Warabiya kids. Definitely not the last.

\begin{center}\rule{0.5\linewidth}{0.5pt}\end{center}

\subsection{Part II, Jam and Other
Escalations}\label{part-ii-jam-and-other-escalations}

Some days they came back carrying stories. Other days, proof. Once, it
was a dead \emph{habu} snake, limp and heavy, held up like a trophy by
hands still small enough to belong in classrooms. Another time, a
massive moray eel, its body coiled in a bucket, teeth still looking
offended. Through it all, the same thing never changed, their smiles.
Calm, proud, completely unshaken, as if this was simply what afternoons
were for.

That's also how I learned Akiko-san kept chickens. Not as a hobby, but
as part of the rhythm of the place. She sold eggs when she had enough,
nothing more, nothing less. The eggs were incredible, deep yellow, rich,
the kind that remind you food has a past before the pan. That's also
when I met her own daughters Renka and Hanaka, moving through the house
with easy confidence, helping without being asked. Everyone seemed oddly
competent here, not in a showy way, but in that quiet, reassuring manner
that makes you realize you're the least prepared person in the room ---
and somehow, that's okay.

After a long trek, I crossed to the neighboring island and found
perfectly ripe adan fruits, heavy, geometric clusters growing straight
out of a tree that looks half-designed by nature, half by stubbornness.
The adan (Pandanus tectorius), sometimes called screw pine, is
everywhere in Okinawa: along beaches, at the edge of forests, standing
firm against salt, wind, and especially \ldots{} typhoons. Its fruit
looks more like an artifact than something edible, a mosaic of hard,
angular segments, bright orange when ripe, stubbornly resistant to the
idea of being food, like a Devil Fruit pulled straight out of One Piece,
mysterious, powerful-looking, and clearly not meant to be eaten without
consequences.. In brief, I obviously needed it ! You wouldn't guess it
at first glance, but with time, heat, and patience, it slowly gives in,
releasing a dense, fibrous flesh with a flavor somewhere between
pineapple and overripe mango, the kind of taste that feels ancient, as
if it fed people long before anyone thought of supermarkets. Since I had
way too much, and baking cookies with only a Japanese fish grill is an
act of pure faith, I brought some to my neighbors.

They were impressed. Curious. \emph{Oishii desu ka?}.. It was
apparently.

\begin{figure}[H]

{\centering \pandocbounded{\includegraphics[keepaspectratio]{.attachments/dd71657358154536b8543860aa0364e78251e921.jpeg}}

}

\caption{image}

\end{figure}%

That evening, Akiko-san invited me in. I hesitated, Japan teaches you to
refuse three times, but eventually accepted. The kids had caught a
couple of \emph{katsuo} (bonito fish). One was cleaning it, another
slicing sashimi, another skewering for tataki. Vegetables grilled over a
\emph{shichirin} (japanese bbq grill). Everyone singing Japanese songs
everyone somehow knew. Calm, coordinated chaos, conducted gently by
Akiko-san.

Watching them move around the house, it slowly sank in how
self-sufficient they were. Not pushed, not supervised every second, just
quietly encouraged to figure things out on their own. They learned by
doing, by failing, adjusting, trying again, and it showed in the way
they handled knives, food, animals, even each other. There was a kind of
calm self-assurance in them that felt rare, something I'd never really
seen in kids that age, and it made me wonder how different growing up
could feel when trust comes before control.

A young guy with long, tied-back hair came over and introduced himself
with a small bow, offered me a beer. Keita-san, a staff helper, gentle
and reserved, already dreaming aloud about building a place like this
one someday.

The food was incredible. Sesame oil in soy sauce with sashimi chef's
kiss! Shintaro, one of the oldest kids explained song lyrics to me. I
introduced him to \emph{Def Tech}. We discussed what a ``yankee'' was.
Cultural exchange complete.

\begin{figure}[H]

{\centering \pandocbounded{\includegraphics[keepaspectratio]{.attachments/377038ea31787eee99b6b447c9f1f66db4a86e1f.jpeg}}

}

\caption{image}

\end{figure}%

As the night wound down, Akiko-san approached me.

``Can I ask you a favor?''

``Mochiron.'' (Of course)

She smiled and walked away.

I had no idea what I'd just agreed to\ldots{}

\subsection{Part III, The Boar}\label{part-iii-the-boar}

A hunter had caught a wild boar in Aharen. Invasive species. They were
bringing it over. They needed help butchering it.

I was not ready for that sentence..

Still, we prepared. I sharpened knives. Plastic sheets went down. Keita
put the kids to bed.

Then Reiji, Akiko-san's husband, arrived in a kei-truck that looked
seconds away from tipping over from the weight.

The boar was enormous. Tusks. Head like a battering ram. We lifted it
together, tied the legs, and Reiji worked quickly and precisely. Calm.
Focused. Blood everywhere. No hesitation --- he moved with a calm
precision I'd only ever seen in a Japanese person.

We skinned it. Renka, the older daughter worked beside me, confidently,
skillfully. No ceremony. Just work.

By nine pm we were deep into it, I was in charge of holding the beast
while the others were cutting it into chunks and bagging them.

Reiji was completely focused, unfazed by the blood splashing across his
glasses or by the raw, overpowering smell of the insides, a smell I was
doing my best to endure without showing it. The stench was intense,
almost unreal, but I was holding myself together.

At one point, he simply handed me the knife and said, ``Now you do.'' I
froze, stiff as a stone. Then I gathered what courage I had left,
grabbed the leg, studied the joints, and went at it. I couldn't stop
thinking about the \emph{jamón serrano} I grew up with in Spain --- a
whole pig leg hanging quietly in the back of the kitchen. This was that
memory stripped of all distance, the real thing.

It made me realize how much we separate animals from food, how far we've
distanced ourselves from the process itself. But these people --- these
kids --- hadn't. They knew exactly what it meant. They didn't flinch,
and because of that, they seemed to respect the animal even more. Not a
single piece was left on the bone; everything that could be eaten was
used.

We eventually wrapped it up with Reiji. An \emph{otsukaresama desu}
(thank you for your hard work) and he was already gone, off to clean the
blood from the boar's head.

I went back inside. After asking how it had been, Akiko-san handed me a
small plastic bag of meat as a token appreciation. I thanked them for
hosting me, and for letting me be part of it, and then headed home. It
was half past midnight. The moon was almost full, bright enough to
follow me the whole way.

When I entered the guesthouse, I startled a few guests who were still
hanging out in the kitchen having beers. Understandably so --- a man
covered in blood, carrying a suspicious bag, isn't exactly reassuring. I
didn't have the energy to explain. I just shrugged, put the meat in the
freezer, and headed for the shower.

I watched the blood mix with the water, spreading and thinning like ink.
I thought about all those crime shows where murderers stare at the
spiral of the sink as the last traces disappear. Strangely, it was
beautiful. Hypnotic, even\ldots{}

\begin{center}\rule{0.5\linewidth}{0.5pt}\end{center}

\subsection{嵐の中の釣り}\label{ux5d50ux306eux4e2dux306eux91e3ux308a}

放課後になると、子どもたちを毎日のように見かけた。走り、自転車に乗り、肩には釣り竿。いつも違う場所へと急いでいく。どうして「いい釣り場」の情報が、あんなにも早く彼らの間を巡るのか、不思議でならなかった。島仕様の、低技術で、子どもたち自身が動かす通信網のようなものが、きっとあったのだろう。

ある午後、彼らは港に集まっていた。その日はたまたま私の休みだった。だから私は、やたらとピンクで、やたらとカワイイ折りたたみ自転車に飛び乗り、釣り竿を掴んで戻り、何事もなかったかのように合流した。

最初、彼らは少し恥ずかしそうだった。ちらりとした笑顔、すぐに糸へ戻る視線。ふざけている子もいたが、多くは真剣そのものだった。本当に大切なことに向き合う時の、あの子ども特有の集中の仕方で。

そこへユウマが現れた。レイさんの息子だ。サニーコーラルで、ベイブレードとポケモンを通じてすでに仲良くなっていた。それだけで友情としては十分だった。彼はまっすぐ私のところへ来て、桟橋釣りがどういうものかを、とても丁寧に説明してくれた。私はスイスの山で川釣りばかりしてきた。この海は、まったく別の世界だった。

彼は私に好きなポケモンを聞いた。私は「ストライク」と言いたかったが、日本語名が分からず、小さく、戦略的な嘘をついた。

「リザードンが好きです。」

彼は勢いよく頷いた。承認完了。すると私が少し日本語を理解できると分かった瞬間、子どもとお年寄りがよくやるように、彼は一気に全速力で話し始めた。言葉の火炎放射器だ。私は同じくらい勢いよく頷き返し、理解の代わりに熱意で乗り切ろうとした。

雨は最初、小雨だった。誰も動かなかった。やがて本降りになっても、誰も動かなかった。そしてついに、横殴りで容赦のない、沖縄らしい雨になった。

ゴアテックスのジャケットを着ていても、私は完全にずぶ濡れだった。だからもう、雨に身を任せることにした。何人かのお母さんたちが子どもを迎えに来た。ユウマに手を振った。それでも少人数の子たちは残り、まったく気にする様子もなく、竿を水に垂らし続けていた。

その時、傘を抱えたアキコさんが現れた。止めるためではない。そんなことをしても無駄だと、よく分かっていたのだ。ただ、この状況を少しでも「耐えられるもの」にするために。背が高く、落ち着いた佇まいで、混沌を整理してしまうような人だった。完璧な英語で、ずいぶん熱心な釣り人ですね、と冗談を言った。

私たちは雨の中で話した。私はゲストハウスを手伝っていることを話した。彼女は、子どもたちは自分の子ではなく、「わらびや」という家で一緒に暮らしていることを教えてくれた。東京や大阪など、全国から来た子どもたち。海と森のそばで、やってみながら学ぶために、親が送り出すのだという。私は、子どもの頃、危うく入れられそうになったフリブールのプロテスタント寄宿学校を思い出した。発想は同じ。実行はまるで違う。

その時、家の近くで見かけていた看板の意味が、ようやく腑に落ちた（わらびや）。

雨の中で三十分ほど話したあと、私たちは家路についた。それが、わらびやの子どもたちと過ごす最初の日だった。最後になるはずがなかった。

\subsection{第二幕 ―
ジャムと、その先へ}\label{ux7b2cux4e8cux5e55-ux30b8ux30e3ux30e0ux3068ux305dux306eux5148ux3078}

ある日は、彼らは話だけを持ち帰ってきた。別の日には、証拠を。

一度は、ぐったりとしたハブ蛇。まだ教室にいそうな小さな手で、戦利品のように掲げられていた。別の日には、巨大なウツボ。バケツの中でとぐろを巻き、歯はまだ怒っているようだった。それでも、変わらないものがあった。彼らの笑顔だ。落ち着いていて、誇らしく、まるでこれが普通の午後だと言わんばかりだった。

その頃、アキコさんが鶏を飼っていることを知った。趣味ではなく、ここでの生活の一部として。卵が十分に集まったら売る。ただそれだけ。卵は本当に素晴らしかった。濃い黄色で、豊かで、食べ物にはフライパンに入る前の「過去」があることを思い出させてくれる味だった。

その時、娘のレンカとハナカにも会った。二人とも自然体で家を行き来し、頼まれずとも手を貸していた。ここでは皆が不思議なほど有能だった。見せびらかすようなものではなく、安心させる静かな確かさ。自分が一番準備不足な人間だと気づかされる。でも、それでいいのだと思えた。

長いトレッキングの後、隣の島へ渡り、完熟したアダンの実を見つけた。自然と頑固さが半分ずつ設計したような木に、幾何学的で重たい房がぶら下がっている。アダン（タコノキ）は沖縄中にあり、浜辺や森の縁で、塩や風、そして台風にも耐えて立っている。

その実は、食べ物というより遺物のようだった。角ばった断片の集合体で、熟すと鮮やかな橙色になる。まるでワンピースの悪魔の実のように、力を秘め、簡単には食べさせてくれなさそうだった。要するに、どうしても試さずにはいられなかった。

見た目からは想像できないが、時間と熱と忍耐をかけると、繊維質で濃厚な果肉を少しずつ現す。味はパイナップルと完熟マンゴーの中間のようで、スーパーが生まれるずっと前から人を養ってきたかのような、古い味がした。

量が多すぎたのと、魚焼きグリルだけでクッキーを焼くのは信仰に近い行為だったので、私は近所にお裾分けした。

彼らは感心し、興味を示した。「おいしいですか？」どうやら、そうだったらしい。

その夜、アキコさんが家に招いてくれた。三回断るのが礼儀だと日本は教えてくれるが、結局、私は受け入れた。

子どもたちはカツオを何匹か釣ってきていた。一人が捌き、一人が刺身を切り、一人がタタキ用に串を打つ。七輪で野菜を焼き、皆が知っている日本の歌を歌う。落ち着いた、統率のとれた混沌を、アキコさんが静かに取り仕切っていた。

彼らが家の中を動き回る様子を見ているうちに、彼らがどれほど自立しているかが、じわじわと伝わってきた。押し付けられず、常に監視されることもなく、ただ自分で考えることを促されている。失敗し、調整し、また試す。その積み重ねが、包丁の扱いにも、食べ物にも、動物にも、互いへの接し方にも表れていた。

その落ち着いた自信は、同じ年頃の子どもたちでは、ほとんど見たことがないものだった。信頼が管理に先立つと、成長はこんなにも違うのかと考えさせられた。

長い髪を後ろで結んだ若い男性が、軽くお辞儀をして自己紹介し、ビールを差し出してくれた。ケイタさん。穏やかで控えめなスタッフで、いつかこんな場所を作りたいと静かに語っていた。

料理は素晴らしかった。刺身に醤油と胡麻油。完璧だ。年長のシンタロウが歌詞を説明してくれた。私はDef
Techを紹介した。「ヤンキー」とは何かを語り合った。文化交流、完了。

夜も更けた頃、アキコさんが近づいてきた。

「お願いしてもいい？」

「もちろん。」

彼女は微笑んで、去っていった。

自分が何に同意したのか、その時はまだ知らなかった。

\subsection{第三幕 ― 猪}\label{ux7b2cux4e09ux5e55-ux732a}

アハレンで猟師が猪を仕留めた。外来種だ。それを運んでくる。解体を手伝ってほしい。

その一文に、私はまったく準備ができていなかった。

それでも、準備は始まった。包丁を研ぎ、ビニールシートを敷き、ケイタは子どもたちを寝かしつけた。

やがて、アキコさんの夫、レイジが軽トラックで到着した。積み荷の重さで、今にもひっくり返りそうだった。

猪は巨大だった。牙、破城槌のような頭。四人で持ち上げ、脚を縛り、レイジは静かで正確な動きで作業を進めた。血は飛び散ったが、迷いはなかった。その落ち着いた正確さは、私が見たことのある中で、日本人にしかできないものだった。

皮を剥いだ。長女のレンカが隣で作業した。自信に満ち、手慣れていた。儀式ではない。ただの仕事だ。

夜九時には佳境に入り、私は獣体を押さえる役目だった。ほかの人たちが肉を切り分け、袋に詰めていく。

レイジは完全に集中していた。眼鏡に血が飛び、内臓の強烈な匂いが漂っても動じない。その匂いに耐えるのが、私には精一杯だったが、なんとか平静を保っていた。

ある時、彼はただ包丁を渡して言った。「次は、君だ。」

私は石のように固まった。それでも残った勇気をかき集め、脚を掴み、関節を見極め、刃を入れた。スペインで育った頃、台所の奥に静かに吊るされていたハモン・セラーノの脚を思い出した。あの記憶から距離をすべて取り除いた、現実そのものだった。

それは、私たちがどれほど動物と食べ物を切り離してしまったかを思い知らされた瞬間だった。しかし、ここにいる人たち、子どもたちは違った。意味を理解していた。ひるまず、だからこそ、動物への敬意がより深かった。骨に残るものは何一つなく、食べられるものはすべて使われた。

作業はようやく終わった。「お疲れさまでした」と言い合い、レイジは猪の頭の血を洗いに行ってしまった。

私は家に戻った。様子を聞かれたあと、アキコさんが小さなビニール袋に入った肉を手渡してくれた。感謝を伝え、参加させてくれたことにも礼を言い、家路についた。

真夜中を少し回った頃だった。ほぼ満月の月が、帰り道をずっと照らしていた。

ゲストハウスに入ると、キッチンでビールを飲んでいた数人の宿泊客を驚かせてしまった。血だらけで、怪しい袋を持った男が現れたら、無理もない。説明する気力はなく、肩をすくめ、肉を冷凍庫に入れ、シャワーへ向かった。

血が水に混じり、インクのように広がって薄れていくのを見ていた。犯罪ドラマで、殺人犯が排水口の渦を見つめる場面を思い出した。不思議なことに、それは美しかった。催眠的ですらあった。

\bookmarksetup{startatroot}

\chapter{}\label{section}

Morino-san looks tough, a little grumpy maybe, the kind of man you might
hesitate to interrupt. Lean and strong, always wearing his sports
sunglasses, he has the look of a professional cyclist. And yet, he is
one of the kindest people I have met.

At heart, all he really wants is to be on the water, or on his boat, a
beautiful fourteen-foot ship, simple and honest, just like him. A true
Osaka man. A man of few words, but always ready for a pun, a quiet joke,
or a sudden burst of absurd humor. One of my fondest memories is when he
unexpectedly imitated the terrifying scream of a Tasmanian devil.

He loves his children deeply. A devoted, protective father, always
attentive without ever being overbearing, present without needing to
show it. His connection to the sea is so natural that he doesn't even
like fishing, nor eating fish. He just likes to see them. I remember him
telling me that very simply, that what mattered was watching them move,
nothing more.

Sometimes I imagine his past life as a truck driver on the Kansai roads,
long hours behind the wheel, his mind drifting elsewhere. I picture him
driving endlessly, dreaming of the Kerama blue, of the Keramashoto sea.
I'm certain that if it were possible, he would have driven his truck
straight through the ocean.

He doesn't speak much. But when he does, it always matters.

Thank you, from the bottom of my heart. You are the one who offered me a
place to stay without really knowing me. You met me once, and you
welcomed me into this island, into your family, into your life.

That gesture alone says everything about the kind of man you are.

Know that I will always see you as a model, someone I look up to. A
great father. A man who works hard for his dream, quietly, steadily. I
respect you deeply, and I respect your family.

Thank you.

\begin{center}\rule{0.5\linewidth}{0.5pt}\end{center}

森野さんは一見厳しく、少し不機嫌そうに見える。声をかけるのを躊躇してしまうようなタイプだ。痩せていて力強く、いつもスポーツサングラスをかけている。どこかプロのサイクリストのような風貌でもある。しかし実際には、彼は私が出会った中で最も優しい人の一人だ。

心の底では、ただ海の上にいたい、あるいは自分のボートの上にいたいだけなのだ。美しい14フィートの船は、彼自身のようにシンプルで誠実だ。まさに大阪の男。口数は少ないが、駄洒落や静かなジョーク、そして時折突然飛び出す荒唐無稽なユーモアには、いつでも応じてくれる。中でも忘れられない思い出のひとつは、彼が突然タスマニアデビルの恐ろしい叫び声を真似した時のことだ。

彼は子どもたちを心から愛している。献身的で守るような父親で、決して押し付けがましくなく、目立たなくても常に気を配っている。海とのつながりはあまりにも自然で、釣りも、魚を食べることさえ好まない。ただ眺めるのが好きなのだ。彼が私にそう言ったのを覚えている。大切なのは、魚が泳ぐのを見ること、それだけだと。

時々、彼の過去を想像する。関西の道を走るトラック運転手として、長い時間ハンドルを握りながら、心は別の場所を漂っていた頃を。果てしなく走り続けながら、ケラマブルー、ケラマ諸島の海を思い描いていたのだろう。もしできるなら、彼はトラックをそのまま海へ走らせていただろう、そう確信している。

彼は多くを語らない。だが、口を開くとき、その言葉にはいつも意味がある。

\bookmarksetup{startatroot}

\chapter{}\label{section-1}

With Miki-san, food moved between us the way words often couldn't. A
dish left on the table for me. A bowl covered carefully, saved for
later. Something warm handed over without explanation, as if it didn't
need one.

We cooked for each other. Not together, but toward one another. She
would prepare something special and place it in front of me, watching
closely, not to judge, just to see if it landed well. I tried to answer
in my own way, Italian meals made with what I had, coffee prepared
carefully, cookies that carried more intention than technique.

Recipes were exchanged loosely. Not written, not measured. Just
described, approximated, trusted. Sometimes something worked. Often it
didn't. It never mattered. What mattered was the gesture, the effort
moving back and forth, the sense that giving was already enough.

In that rhythm of offerings, something settled between us. No big
declarations. Just attention, care, and the quiet understanding that
hospitality doesn't need to be explained when it's sincere.

Thank you, Miki-san, for welcoming me so fully, without conditions or
expectations. Thank you for the meals, for the patience, for the trust,
and for the space you gave me to be myself in your home. You made me
feel seen, cared for, and included in ways that go far beyond words. I
will always carry this with me, and I am deeply grateful to call you not
only my host, but my friend.

\begin{center}\rule{0.5\linewidth}{0.5pt}\end{center}

美樹さんとの間では、言葉では伝えきれないことが、食べ物を通して自然と行き交っていた。私のためにテーブルに残された一皿。後で食べられるよう、丁寧に蓋をされた茶碗。理由も説明もなく、そっと手渡される温かいもの。まるで、説明など最初から必要なかったかのように。

私たちは互いのために料理をした。一緒に、ではなく、互いに向かって。彼女は特別な一品を用意し、私の前に置くと、静かに様子を見ていた。評価するためではなく、ただ受け取ってもらえるかを確かめるために。私は私なりの方法で応えようとした。手元にある材料で作るイタリア料理、心を込めて淹れたコーヒー、技術よりも気持ちを重ねたクッキー。

レシピはゆるやかに交わされた。書き留めることも、正確に量ることもなく、ただ言葉で伝え、感覚で補い、互いを信じた。うまくいくこともあれば、そうでないことも多かった。それでも問題はなかった。大切だったのは結果ではなく、その行為そのもの、行き交う気持ち、与えること自体がすでに十分だったという感覚だった。

そうした贈り合いのリズムの中で、私たちの間には静かに何かが落ち着いていった。大げさな言葉は必要なかった。ただ注意深さと気遣い、そして誠実なもてなしには説明などいらない、という穏やかな理解がそこにあった。

\bookmarksetup{startatroot}

\chapter{}\label{section-2}

Kobayashi-san has a presence that stays with you. A slender silhouette,
always slightly in motion, a relaxed posture shaped by years spent
outdoors rather than indoors. He smiles easily, almost constantly, with
that gentle dad humor that arrives quietly and lands perfectly. Nothing
forced, nothing heavy, just a natural ease, as if curiosity itself were
a habit he never dropped.

In the kitchen, he moves the same way he does everywhere else: with
focus, generosity, and an obvious pleasure in doing things properly. He
showed me the ropes of traditional Japanese cooking, sukiyaki, tonkatsu,
ozoni, not as recipes to memorize, but as gestures to feel. Timing over
measurement. Heat over theory. He cooks the way he lives: hands-on,
precise, fully engaged. His passion spills into food effortlessly, as if
caring deeply were simply his default setting.

That same precision carries into the way he observes the world. He
studied ornithology, and it shows. His attention is methodical, patient,
grounded. He doesn't speculate or embellish. He reads nature carefully,
on its own terms.

In his twenties, he went to the wilds of the Yukon. No elaborate plan,
no preparation worthy of the name. Canoe, wilderness, months on the
river. A true leap into the unknown, guided more by instinct than by
strategy. He tells it casually, almost offhandedly, as if drifting
through the Canadian wild without a shared language were just another
reasonable option life presented him with.

He didn't plan a life. He followed a current.

Today, he lives on a tiny island, far from the rivers and mountains that
shaped him. Surrounded by salt water --- which he admits, somewhat
ironically, he doesn't particularly like --- yet perfectly content to be
immersed in nature all the same. The setting changed, but the movement
didn't. He is still in motion.

That motion defines him. Always doing something. Cutting grass under the
scorching sun. Renovating houses. Cycling up and down the island to
help, fix, carry, adjust. He never waits to be asked. He never
complains. Movement, for him, isn't restlessness --- it's balance.
Staying alive by staying engaged.

He embodies \emph{ikigai}. Not as an idea, but as a practice. No desk
for forty years. No abstract philosophy. Just work that matters, effort
directed somewhere real, days that leave a trace.

There's something deeply striking about him: a youthful curiosity
carried in a body shaped by time, without denial, without performance.
He isn't nostalgic. He's moving forward.

Meeting Kobayashi-san feels like a quiet lesson in how to live, not
explained, not taught, simply shown, every day, through action.

\begin{center}\rule{0.5\linewidth}{0.5pt}\end{center}

小林さんは忘れがたい存在感を持っている。細身のシルエットはいつもわずかに動いていて、屋内よりも屋外で過ごしてきた年月がつくった、力の抜けた佇まい。彼は気負うことなく、ほとんどいつも微笑んでいる。静かに届き、ちょうどよく効く父親のようなユーモアを携えて。無理も重さもない。ただ自然体で、まるで好奇心そのものが手放されることのなかった習慣であるかのようだ。

厨房でも、その動きは変わらない。集中力と寛大さ、そして物事をきちんとやることへの純粋な喜び。彼は私に、すき焼き、とんかつ、お惣菜といった日本の家庭料理を教えてくれた。暗記するレシピとしてではなく、身体で覚える所作として。計量よりもタイミング、理論よりも火加減。彼は生き方そのもののように料理をする。手を動かし、正確に、全身で向き合う。深く向き合うことが、彼にとってはごく自然なことなのだ。

その正確さは、世界を観る眼差しにも表れている。鳥類学を学んだ人らしく、観察は系統立っていて、忍耐強く、地に足がついている。憶測も誇張もない。自然を、自然のままに、丁寧に読み取っていく。

二十代の頃、彼はユーコンの大自然へと向かった。綿密な計画も、名に値する準備もないままに。カヌーと荒野、そして川で過ごした数か月。本能に導かれた、真の未知への跳躍だった。彼はそれを何気なく語る。共通の言葉も持たずにカナダの荒野を漂うことが、人生が差し出した選択肢のひとつに過ぎなかったかのように。

彼は人生を計画しなかった。ただ流れに従ったのだ。

今、彼は小さな島に暮らしている。彼を形づくった川や山々からは遠く離れ、塩水に囲まれた場所だ。「あまり好きじゃないんだけどね」と、少し皮肉を込めて彼は言う。それでも自然の中に身を置くことに、彼は満足している。環境は変わっても、動きは変わらない。彼はいまも動き続けている。

その動きこそが、彼を定義している。灼熱の太陽の下で草を刈り、家を直し、自転車で島を行き来しながら、助け、修理し、運び、整える。頼まれるのを待つことはない。文句も言わない。彼にとって動くことは落ち着きのなさではなく、均衡なのだ。関わり続けることで、生き続けている。

彼は生きがいを体現している。理念としてではなく、実践として。四十年間机に向かう人生でもなく、抽象的な哲学でもない。ただ意味のある仕事、現実に向けられた努力、痕跡を残す日々がある。

彼には胸を打つ何かがある。時を重ねた身体に宿る、若々しい好奇心。否定も虚飾もない。彼は懐かしむ人ではない。前へ進み続けている。

小林さんとの出会いは、生き方についての静かな教訓のようだ。説明されることも、教えられることもない。ただ、日々の行動を通して示されるだけなのだ。

\bookmarksetup{startatroot}

\chapter{}\label{section-3}

The shrill rasp of cicadas, just awakened from their long underground
hibernation, is deafening. One after another they scream, weaving a
shattering unison, louder than the sum of each individual clicking
joined together.

\begin{figure}[H]

{\centering \pandocbounded{\includegraphics[keepaspectratio]{.attachments/9243b465682158168dd76cacd7899090755d9f91.jpeg}}

}

\caption{image}

\end{figure}%

I then let myself be gently rocked in a fabric hammock, suspended, or at
least I like to believe so, not by this rope that has clearly felt the
weight of passing years, but by the benevolent spirit of the \emph{kami}
from the nearby temple.

\begin{figure}[H]

{\centering \pandocbounded{\includegraphics[keepaspectratio]{.attachments/60ebdb56d98ad507e464a240128406b09ea1ed4c.jpeg}}

}

\caption{image}

\end{figure}%

A breath of sea air draws me back to observing my surroundings. I look
at my aluminum container that serves as my home. I've learned to grow
attached to it, but more than that, to the life it offers. Dew sparkles
on the thin, rain-soaked \emph{adan}. Farther away, crows peck at fallen
\emph{shikuwasa}, small sharp citrus fruits that reveal their full
flavor in a glass of Okinawan awamori.

\begin{figure}[H]

{\centering \pandocbounded{\includegraphics[keepaspectratio]{.attachments/d432ab71bd23d09fd107ff1783525338274bc9cb.jpeg}}

}

\caption{image}

\end{figure}%

\begin{center}\rule{0.5\linewidth}{0.5pt}\end{center}

長い地中での眠りから目覚めたばかりの蝉たちの、甲高くざらついた鳴き声が耳を打つ。ひとつ、またひとつと鳴き始め、それぞれの音が重なり合い、個々の響きの総和をはるかに超える、砕けるような大きなうねりをつくっていく。

やがて私は、布のハンモックに身を預け、ゆっくりと揺られる。吊るしているのは、年季の入ったこの縄ではなく、近くの神社に宿るやさしい神々なのだと、そう思うことにしている。

ふと、潮の香りを含んだ風が頬をなで、意識を現実へと引き戻す。視線を巡らせると、住まい代わりのアルミのコンテナが目に入る。いつの間にか、ただ慣れただけでなく、そこで営まれる暮らしそのものに愛着を覚えるようになっていた。

雨に濡れた細いアダンの葉には、露がきらきらと光っている。少し離れたところでは、カラスたちが落ちたシークヮーサーをついばんでいる。その小さく鋭い柑橘は、沖縄の泡盛に落としたとき、ようやく本来の味をほどく。

\bookmarksetup{startatroot}

\chapter{Bon Voyage Chef !}\label{bon-voyage-chef}

Some days, this island feels like a dream, and nothing embodies that
more than Mai-san's food truck : ``Bon Voyage''.

She's in her twenties tell but quite fired up, and she chose to trade
the quiet predictability of Okayama for the wild freedom of Tokashiki,
following a dream only she could see here.

\begin{figure}[H]

{\centering \pandocbounded{\includegraphics[keepaspectratio]{.attachments/b1cfe05072269b5b89e66aabb7a420c01d9c1e6a.jpg}}

}

\caption{image}

\end{figure}%

She opens when she feels like it no schedule just instinct. Maybe that's
why the sides of her truck always feels like a mirage in the heat.

Each time I walk toward the beach, I catch myself scanning the curve of
the road, wondering if today is one of these lucky days if the marriage
will turn out to be real this time, and if I will get the chance to
order one of her legendary ``Cuba sando'', sandwich she presses with a
kind of quiet focus that makes you stop whatever you were doing just to
watch the cheese melt out of it.

\begin{figure}[H]

{\centering \pandocbounded{\includegraphics[keepaspectratio]{.attachments/5b27839b8addc25b80d33ab0bb3509699eee596e.jpg}}

}

\caption{image}

\end{figure}%

On this tiny island, there is something magical about not knowing. And
then when Mai-san's stand is open, it feels like the day itself decided
to be kind..

\begin{center}\rule{0.5\linewidth}{0.5pt}\end{center}

ある日は、この島が夢のように感じられる。そして、その感覚を最も体現しているのが、舞さんのフードトラック「ボン・ボヤージュ」だ。

彼女はまだ二十代だが、静かな情熱を内に秘めている。岡山の穏やかで予測可能な日々を手放し、トカシキの奔放な自由を選んだ。ここでしか見えなかった、自分だけの夢を追いかけるために。

営業は気まぐれで、決まったスケジュールはない。ただ感覚に従うだけ。だからこそ、彼女のトラックは炎天下の中で、いつも蜃気楼のように現れるのかもしれない。

浜辺へ向かうたび、私は無意識に道のカーブへ視線を送ってしまう。今日は運のいい日だろうか。その蜃気楼は本物になるだろうか。そして、彼女の伝説的な「キューバサンド」を注文できるだろうか。静かな集中力でプレスされるそのサンドイッチは、溶け出すチーズを眺めているだけで、思わずすべての動きを止めてしまう。

この小さな島では、わからないままでいること自体に魔法がある。そして舞さんの屋台が開いている日は、まるでその日そのものが、優しくなろうと決めたかのように感じられる。

\bookmarksetup{startatroot}

\chapter{}\label{section-4}

Like at school, everything here follows the sound of soft, and familiar,
alarms. You don't need a watch --- just open your ears. At 8 a.m., the
daily announcement lets everyone know what's happening and whether the
boat will manage to cross the swell. At 12:30, the midday message brings
news of the weather and the day's activities.

At 5 p.m., the first chime rings out, warning children that it's time to
head back home. Then at 5:30, the final bell, \emph{going home}, sounds.
It gives a strangely formal air to the sunset, as if telling the night
stars their break is over and it's time to get back to work.

And of course, it signals to the children that playtime is done. The
island shifts. Space opens for the snakes, the centipedes, the wild
boars, and the other creatures of the night.

\begin{center}\rule{0.5\linewidth}{0.5pt}\end{center}

学校のように、ここではすべてが柔らかく、どこか懐かしい音の合図に従って進んでいく。時計はいらない。ただ耳を澄ませばいい。朝8時、その日の予定や、波の具合で船が出せるかどうかが放送で知らされる。12時半には、天気やその日の動きについての昼の知らせが流れる。

夕方5時、最初のチャイムが鳴り、子どもたちに帰る時間が近いことを告げる。そして5時半、最後の合図------「帰宅」の放送が響く。夕暮れに、どこか儀式めいた空気が差し込む。まるで夜の星たちに、休みは終わりだ、そろそろ持ち場に戻る時間だと伝えているかのように。

もちろんそれは、子どもたちにとっても遊びの終わりの合図だ。島の気配が少し変わる。空間がゆるみ、ヘビやムカデ、イノシシ、そして夜の生き物たちのための時間がひらけていく。

\bookmarksetup{startatroot}

\chapter{Jazz \& Kare}\label{jazz-kare}

Inside Sunny Coral time moved differently Rai-san is steering his
Japanese curry in the kitchen behind, his jazz playlist drifting to the
café, warm and grainy the kind of song that wraps around you before you
even notice that. I sat there scratching his cat's neck (Kotora 子虎 -
Small tiger) . Eating this kare, felt like what jazz must tastes like.

Deep notes, bright notes, little surprises tucked between bytes.
Outside, the island's wind shimmered, but inside this small hut, it was
all mellow rhythm and gentle conversations. For a moment, I wasn't just
traveling, I was drifting, letting a bowl of curry and a few jazz chords
tell me everything I needed to know about being alive and far from home
in the best possible way.

\begin{center}\rule{0.5\linewidth}{0.5pt}\end{center}

サニーコーラルの中では、時間の流れが少し違っていた。奥のキッチンでは、レイさんが日本のカレーを静かにかき混ぜている。その間、ジャズのプレイリストがカフェ全体に漂ってくる。温かく、少しざらついた音色。気づく前に、そっと体を包み込んでくるような曲だ。

私は席に座り、彼の猫、子虎
の首元を撫でていた。このカレーを食べていると、ジャズがどんな味なのか、少しわかる気がした。

深い音、明るい音。一口ごとに隠された小さな驚き。外では島の風がきらめいていたけれど、この小さな小屋の中では、すべてがまろやかなリズムと穏やかな会話に包まれていた。

ほんのひととき、私はただ旅をしているのではなく、漂っていた。一皿のカレーと、いくつかのジャズのフレーズが、生きていること、そして故郷から遠く離れながらも、最良のかたちでここに在るということについて、知るべきすべてを静かに教えてくれた。

\bookmarksetup{startatroot}

\chapter{The sea}\label{the-sea}

It lies there, supple and unchanging. Around me, white birds and rusted
houses. In the breaking of the waves, there is something of your voice.
In the foam, a glint of your smile.

The sea rocks my soul the way it did that evening, when a song by
Charles Trenet brought us closer, a little more with each chorus. Now,
every wave carries me back to that moment.

\begin{center}\rule{0.5\linewidth}{0.5pt}\end{center}

\subsection{海}\label{ux6d77}

そこにある。しなやかで、変わらない。周りには白い鳥と、錆びた家々。波が砕ける音の中に、どこかあなたの声が混じっている。泡のきらめきに、あなたの笑顔がふとよぎる。

あの夕暮れと同じように、海は心を揺らす。シャルル・トレネの一曲が流れ、サビのたびに少しずつ距離が近づいていったあの時間。いまでは、一つひとつの波が、あの瞬間へと静かに連れ戻していく。

\bookmarksetup{startatroot}

\chapter{The Smell of Wet Towels}\label{the-smell-of-wet-towels}

One of those days, could've been the day before\ldots{} but no. Pretty
sure it was the next one..

I was biking through the village, just trying to get some fresh air,
when I started to smell it. Something familiar. Something I knew very
well.

But I couldn't quite place it. I found myself sniffing the air like a
wild animal, short breaths, searching, nostrils flaring like antennas.
The smell was everywhere.

And then it hit me.

Wet towels. But not just any wet towels.

A towel after a long day at the beach. Left in a bag. Still damp,
starting to turn. That specific, heavy, salty smell , right on the edge
of becoming something else.

It was as if the sea itself had packed a bag and left it open in the
village, letting its scent drift through the streets for no reason at
all.

It brought me straight back to the Cantabrian Sea when I was a child.
Long days in the water with my cousins, fighting the waves like it meant
something. The Atlantic was always freezing, even in summer. We'd run
back to the pool every ten minutes just to warm up.

Our lips would turn purple, but we'd go straight back in anyway.

And somehow, that smell, that damp towel smell, had been waiting there
all along. Decades later, it found me again.

\begin{center}\rule{0.5\linewidth}{0.5pt}\end{center}

そんな日のひとつ。前の日だったかもしれない\ldots\ldots いや、違う。たぶんその次の日だ。

村の中を自転車で流していた。少し空気を吸いたくて。そんなとき、ふと匂いに気づいた。どこか懐かしい、よく知っている匂い。

でも、うまく思い出せない。私はまるで野生の動物みたいに、短く何度も息を吸い込んで、空気を嗅ぎ回っていた。鼻をアンテナみたいに広げて。匂いはどこにでもあった。

そして、ふっとわかった。

濡れたタオルの匂い。でも、ただの濡れたタオルじゃない。

一日中ビーチで過ごしたあとのタオル。バッグの中に入れっぱなしで、まだ湿っていて、少しずつ変わり始めている。あの独特の、重たくて、塩っぽい匂い。何か別のものになりかけている、その手前の匂い。

まるで海そのものがバッグをひとつ用意して、それを村のどこかに開けっぱなしにしていったみたいだった。ただ気まぐれに、その匂いを通りに漂わせているみたいに。

それは一瞬で、子どもの頃のカンタブリア海へと私を連れ戻した。いとこたちと一日中海に入って、波と本気で戦っていたあの時間。大西洋の水は、夏でもいつだって冷たかった。私たちは十分おきにプールに戻って、体を温めていた。

唇は紫になっても、またすぐ海に戻っていった。

そして気づけば、その匂い,あの湿ったタオルの匂いは、ずっとそこにあったのだ。何十年も経ってから、また私を見つけに来た。

\bookmarksetup{startatroot}

\chapter{Masturi Season}\label{masturi-season}

\subsection{Tanabata}\label{tanabata}

It started, like many things on the island, quietly. I didn't notice it
at first. Just small changes. People cleaning, moving with a certain
intention. Graves being tended to. Paths cleared. Not in a symbolic way,
but in a very real one, as if someone was expected to arrive, and you
didn't want them to trip on the way in.

Later, Miki-san explained it to me. A story about two lovers separated
by a river of stars. They were only allowed to meet once a year, if they
had both done their part, if they had worked enough, if the sky allowed
it.

A bridge of birds would form.

Or not.

It stayed with me longer than I expected.

Tanabata didn't feel romantic to me. It felt conditional. It made me
think about my past relationship, in which I consistently lost my way
onto it. Which is ironically why I decided to flee and end up here.

\subsection{Obon}\label{obon}

Then came Obon. The island changed rhythm. Not louder, not busier, just
fuller. As if something invisible had taken up space between people.

Families gathered. Food appeared everywhere. Doors stayed open.

The idea was simple: the dead come back for a few days.

No one questioned it. No one dramatized it either. It wasn't mystical.
It was practical.

They come. You welcome them. You spend time together. Then you help them
leave.

I had seen something similar before, in Peru, during Día de los Muertos.
Different colors, different smells, different sounds, but the same
intention underneath.

The same refusal to let memory dissolve.

Memory, given structure.

And then, at night, you would hear it.

\subsection{Eisa}\label{eisa}

At first, just distant drums. Then closer. Then suddenly, they were
there, moving through the streets in a slow, deliberate procession.

Drums strapped to bodies. Movements both sharp and fluid. Voices rising,
falling, repeating.

They weren't performing. They were guiding.

Someone told me the dances had once disappeared. Lost after the war. And
then, somehow, brought back. Rebuilt. Relearned. Not exactly the same,
but alive again.

There were roles within it. The ones carrying the rhythm deep in their
shoulders. The ones leading the song with the sanshin. And the ones with
painted faces, half clowns, half conductors, quietly holding everything
together without ever taking the center.

They stopped in front of houses. In front of newly opened places.

Always moving. Always drumming. Always dancing.

As if the sound itself was a path.

\subsection{The tug-of-war - Oonahiki}\label{the-tug-of-war---oonahiki}

And then, weeks later, everything shifted.

From the past to the future.

One moment I was riding through the narrow, hilly roads of the village,
the next I was being pulled into something I didn't fully understand.
Miki-san brought to the fields to meet with the other women of the
village.

Rice straw, laid out in piles. Dry, golden, deceptively light. We began
twisting it, braiding it, tightening it. Over and over. Small strands
becoming larger ones, larger ones becoming something thicker, heavier,
patterns repeating inside patterns, like fractals you could hold in your
hands.

This wasn't just a rope.

It was \emph{the} Ufunaa.

Two sides of the island working in parallel. Not against each other, not
yet. The East and the West. Fishermen and farmers. Sea and land, in a
quiet kind of balance.

The rope itself came in two parts. One with a loop. The other with a
thick, straight head. Later, they would be joined, locked together with
a massive wooden pin. People laughed about it, of course. No one missed
the implication.

But before that, there was work. Hours of it.

Hands moving without much talking. Older hands correcting younger ones.
A rhythm that didn't need explaining.

At that point, I thought I was just there to watch.

I was wrong.

Morino-san shouted at me. Sharp, direct. No room for hesitation. A
tenugui was shoved into my hands, the mark of the West, tied tight
across my forehead before I could even react.

Then I was in.

Forty people, maybe more. The rope was impossibly heavy. It didn't feel
like something meant to be moved.

We carried it to the center of the village.

The two sides met. The ropes were joined. The wooden pin dropped
through. A moment of pause.

Then tension.

Across from me stood the East. At the front, a man they called Jumbo;
easily the largest man I had seen in Japan, built like the kind of
weight that doesn't get moved, only resisted.

Behind him, Michael, the only other foreigner on the island. An English
teacher from the US. We had spent weeks sharing quiet routines, small
conversations, the kind that build familiarity without effort.

Now we were on rivals. He looked at me and smiled.

``Ready?'' Not really. But it didn't matter. The signal came.

At first, nothing moved. Then a shift. The East pulled hard. We gave
ground. Just a little.

Then something changed behind me. A sound. A rhythm. Voices rising
together.

``Ousss!''

We pulled\ldots{} Not as individuals, but as something else entirely.

I had never pulled that hard in my life.

And suddenly, it was over.

The rope moved. The East collapsed forward.

Afterwards, everything softened again. Laughter. Hands on shoulders.
Beers appearing from nowhere. Stories from previous years, told like
they had just happened.

Everyone ``otsukaresamed'' everyone else.

We dropped into that resting squat, heels flat on the ground, bodies
folded low, close to the earth, breathing slowly again. The same people
who had been straining seconds before were now completely at ease.

The future; a good harvest, which I didn't know I would actually be part
of..

\begin{center}\rule{0.5\linewidth}{0.5pt}\end{center}

\subsection{七夕}\label{ux4e03ux5915}

それは、この島の多くのことと同じように、静かに始まった。
最初は気づかなかった。ほんの小さな変化だけ。人々は掃除をし、どこか目的を持ったように動いていた。お墓の手入れがされ、小道が整えられていた。それは象徴的なものではなく、とても現実的な理由からだった。まるで誰かが来ることになっていて、その人がつまずかないようにしているかのように。

あとで、ミキさんが教えてくれた。星の川に隔てられた恋人たちの話。二人は一年に一度だけ会うことが許される。ただし、それぞれがやるべきことを果たし、十分に働き、そして空がそれを許したときだけ。

鳥たちが橋を作る。

あるいは、作られないこともある。

思っていたよりも、ずっと心に残った。

七夕は、私にはロマンチックには感じられなかった。どちらかというと、条件付きのものに思えた。かつての自分の関係を思い出した。私はその中で、何度も道を見失っていた。そして皮肉なことに、それこそが、すべてから逃げるようにここへ来た理由でもあった。

\begin{center}\rule{0.5\linewidth}{0.5pt}\end{center}

\subsection{お盆}\label{ux304aux76c6}

やがて、お盆がやってきた。島のリズムが変わった。騒がしくなるわけでも、忙しくなるわけでもない。ただ、何かが満ちていくような感覚。目に見えないものが、人と人のあいだに入り込んできたようだった。

家族が集まり、食べ物があふれ、家の扉は開け放たれたままになる。

考え方はシンプルだった。死者が数日間、戻ってくる。

誰もそれを疑わないし、大げさにも扱わない。不思議なことでもない。むしろ実務的だった。

帰ってくる。迎える。一緒に過ごす。そして、また送り出す。

似たようなものを、私はペルーで見たことがあった。死者の日のときに。色も匂いも音も違うけれど、その奥にあるものは同じだった。

記憶を消させないという、静かな意志。

形を与えられた記憶。

そして夜になると、それが聞こえてくる。

\subsection{エイサー}\label{ux30a8ux30a4ux30b5ux30fc}

最初は遠くの太鼓の音。やがて近づき、気づけばすぐそこにいた。ゆっくりと、しかし確かな足取りで通りを進んでいく。

太鼓を体に結びつけ、鋭くも滑らかな動き。声は上がり、下がり、繰り返される。

彼らは演じているのではなかった。導いていた。

戦後、一度は失われた踊りだと聞いた。それが、いつの間にか取り戻され、作り直され、学び直された。まったく同じではないかもしれないが、確かに生きていた。

そこには役割があった。肩にリズムを担う者。三線で旋律を導く者。そして、顔を塗った者たち。道化のようでいて、全体を静かに支える存在。中心に立つことはないが、欠かせない存在だった。

彼らは家の前で立ち止まり、新しくできた場所の前でも足を止める。

常に動き、常に打ち鳴らし、常に踊る。

まるで音そのものが、道になっているかのように。

\subsection{大綱引き（オオナヒキ）}\label{ux5927ux7db1ux5f15ux304dux30aaux30aaux30caux30d2ux30ad}

そして数週間後、すべてが変わった。

過去から未来へ。

ある瞬間まで、私は村の細くて起伏のある道を自転車で走っていた。次の瞬間には、自分でもよくわからない何かに引き込まれていた。ミキさんに連れられて、畑へ行き、村の女性たちと合流した。

そこには稲わらが広げられていた。乾いていて、金色で、見た目よりも軽い。それをねじり、編み、締めていく。何度も何度も。細い束が太くなり、太いものがさらに大きくなる。同じ形が繰り返されていく。まるで手の中にあるフラクタルのようだった。

それはただの縄ではなかった。

「ウフナー」だった。

島の二つの側が並行して作業する。まだ対立ではない。東と西。漁と農。海と陸。そのあいだにある、静かな均衡。

縄は二つに分かれていた。片方は輪、もう片方は太くまっすぐな先端。やがてそれらは一つに結ばれ、大きな木の杭で固定される。その様子に、誰もが笑っていた。意味は、誰の目にも明らかだった。

だが、その前にあるのは作業だった。

何時間も続く作業。

ほとんど言葉はなく、年長者の手が若い手を正していく。説明のいらないリズム。

そのときの私は、ただ見ているだけだと思っていた。

違った。

森野さんが鋭く声を上げた。ためらう余地はなかった。手ぬぐいが手渡され、西の印として額にきつく結ばれた。

気づけば、その中にいた。

四十人ほど。あるいはそれ以上。縄は信じられないほど重く、本来動かすものではないように感じられた。

それを担いで、村の中心へ運ぶ。

東と西が向き合う。縄が結ばれ、木の杭が打ち込まれる。一瞬の静止。

そして緊張。

向こう側には東の人々。その先頭には「ジャンボ」と呼ばれる男がいた。日本で見たことがないほどの大きさ。動かされる存在ではなく、むしろ動きを止める側のような体だった。

その後ろにマイケルがいた。島で出会った、もう一人の外国人。アメリカから来た英語教師。特別なことをしたわけではないが、静かな日常を共有してきた。

その彼が、今は向こう側にいる。

目が合い、彼は笑った。

「準備はいい？」

まったくできていなかった。

だが関係なかった。

合図が出る。

最初は何も動かない。次に、わずかな揺れ。東が強く引く。少しだけ後退する。

そのとき、背後で何かが変わった。

音。リズム。声が一つになる。

「オウス！」

引いた。

個人としてではなく、別の何かとして。

あれほど力を込めたことはなかった。

そして、あっけなく終わった。

縄が動き、東が前に崩れた。

その後、すべてはまた緩んでいく。笑い声。肩を叩き合う手。どこからともなく現れるビール。ついさっきのことのように語られる昔話。

人々は互いに「お疲れさま」と声をかけ合う。

やがて皆、あのしゃがみ方で休む。かかとは地面についたまま、体を低く落とし、静かに呼吸を整える。ほんの少し前まで全力で引いていた人たちが、もう完全に落ち着いている。

未来は決まった。

あるいは、少しだけ良い方向へ動いた。

豊作になると言われていた。

その中に、自分も含まれているとは、そのときはまだ知らなかった。

\bookmarksetup{startatroot}

\chapter{}\label{section-5}

\# Akajima

Ici on raconte l'histoire d'un chien nommé Shiro, qui nageait chaque
jour jusqu'à l'île voisine pour retrouver Marilyn, rencontrée avant que
son maître ne déménage. Trois kilomètres de courants entre Akajima et
Zamami, parcourus inlassablement. Aujourd'hui, deux statues se font face
de part et d'autre de l'océan, rappel silencieux que certaines distances
ne comptent pas quand le lien est juste.

\begin{center}\rule{0.5\linewidth}{0.5pt}\end{center}

\# 阿嘉島

ここでは、シロという犬の話が語り継がれている。飼い主の引っ越しで離れ離れになる前に出会ったマリリンに会うため、毎日泳いで隣の島へ渡っていたという。阿嘉島と座間味島のあいだ、潮の流れが強い三キロの海を、何度も、変わらずに。

いまは海を挟んで向かい合うように二つの像が立っている。ただ静かに、たしかなつながりがあれば、距離はそれほど問題ではないことを示している。

\bookmarksetup{startatroot}

\chapter{Between stone and memory}\label{between-stone-and-memory}

At the far western edge of Japan lies Yonaguni, where the gaze almost
reaches Taiwan, and the sea already feels like it opens onto other
worlds.

Here, facing the underwater formations, the boundary between geology and
architecture begins to blur. Beneath the surface, immense angular
structures rest in silence, terraces, steps, sharp walls, as if the
stone had been cut, ordered, conceived. Natural formation or human
intervention, the debate remains open, and it is precisely there that
the mystery resides.

Yonaguni offers no inscriptions, no certainties, only shapes too regular
to be ignored, too ancient to be easily explained. They sleep beneath
the water, battered by currents, observed by divers the way one
contemplates an unsolved riddle, with equal parts caution and
fascination.

Around these submerged stones drift distant echoes, Tiwanaku in the
Andes, the Maya, the Minoans, the Jōmon culture, Rapa Nui. Civilizations
separated by oceans, yet sometimes connected by similar intuitions,
monumental stonework, cosmologies, relationships to time and to the
earth.

And once again, my grandfather Antonio seems to continue walking beside
me. Not as a fixed memory, but as a diffuse presence, emerging through
landscapes, ancient stones, and the names whispered while facing the
sea. Here, as elsewhere, I feel his gaze resting on these enigmas, he
who believed in lost civilizations and invisible links between worlds.
As if the journey never fully belonged to me, but quietly extended his
own.

Summits as vestiges of a larger continent, perhaps submerged twelve
thousand years ago. Coincidences, or fragmented memory of an ancient
world. Here, the sea does not erase everything, it also keeps its
silence.

\begin{center}\rule{0.5\linewidth}{0.5pt}\end{center}

\bookmarksetup{startatroot}

\chapter{石と記憶のあいだ}\label{ux77f3ux3068ux8a18ux61b6ux306eux3042ux3044ux3060}

日本の最西端、与那国。視線の先には台湾がかすみ、海はすでに別の世界へとひらけているように感じられる。

ここでは、海中の地形を前にして、地質と建造物の境界がゆるやかにほどけていく。水面の下、角ばった巨大な構造が静かに横たわる。段差や階段、鋭い壁面------まるで石が切り出され、整えられ、意図されてそこにあるかのように。自然の産物か、人の手によるものか。その議論は今も決着を見ず、まさにそのあいだに、この場所の謎がある。

与那国は、碑文も確証も残していない。ただ、見過ごすにはあまりに整いすぎ、説明するにはあまりに古い形だけがある。潮に打たれながら海中に眠り、ダイバーたちはそれを、解かれない問いに向き合うように、慎重さと魅了のあいだで見つめている。

その沈んだ石のまわりには、遠い響きが重なる。アンデスのティワナク、マヤ、ミノア、縄文、ラパ・ヌイ。海を隔てた文明でありながら、ときに似た直感を共有しているかのように、巨石の構築や宇宙観、時間や大地との関係が静かに呼応する。

そしてまた、祖父アントニオが隣を歩き続けているように感じられる。固定された記憶ではなく、風景や古い石、海に向かってつぶやかれる名前の中から、ふと立ち現れる気配として。ここでも、ほかの場所と同じように、彼の視線がこの謎の上にそっと置かれているのを感じる。失われた文明や、世界のあいだにある見えないつながりを信じていた彼。その旅は、完全に自分のものではなく、静かに彼の歩みを引き継いでいるかのように。

山頂は、かつての大陸の名残かもしれない。およそ一万二千年前に沈んだとも言われる大地の断片。偶然なのか、それとも古い世界の断片的な記憶なのか。ここでは、海はすべてを消し去るのではなく、その沈黙ごと抱え込んでいる。

\bookmarksetup{startatroot}

\chapter{Sea Rain.}\label{sea-rain.}

Mimitsu-machi has been left untouched since 1830, and it's the place
where the first emperor of Japan, Jimmu, 2600 years ago, set sail from
there to unite Japan. I just came there because it was close to the most
famous surf breaks of Jpapan. But I was fortunate enough to meet this
guy, Keizo.

A middle-aged surfer, retired, he's guided most surfers from all around
the world. His English was a little rusted, and it was sometimes
difficult to even understand his Google translations, but he was the
greatest Japanese person I've met so far. He was burned by the sun,
which is very uncommon to see in Japan, as they are all very cautious
about the sun.

Plus, having been on many foreigners for several years, you can see that
his Japanese manners had certainly been corroded by Western surfer life.
He explained to me how he got knocked out from an 6-meter storm wave,
which dislocated both of his knees three years ago, and that he hasn't
been able to surf since then. But that this weekend, he was going to
surf again for the first time in three years with me.

I felt very honored, but I was also scared that he would bring me to a
14 feet wave break at the same time. Anyhow, it was my first day, and my
face was battered by the shores of the Miyazaki prefecture. Suffice to
say, the sea wasn't calm. There I was, waiting for the waves to jump on
me and my board. I was just amazed by the sheer size of the break.

You could count over 50 surfers riding along, and that was a big thing.
The sun had just set on the sea about half an hour ago, but the clouds
were rolling in. They were robbing us of a clear sense of direction.
They were flashing like a megawatt each minute or so. Instead, we were
soon welcoming the morning sea rain. What a feeling.. Sensing the
droplets leaping onto your face, while you're paddling out in the sea
rain.

Thanks Keizo !

\begin{center}\rule{0.5\linewidth}{0.5pt}\end{center}

\bookmarksetup{startatroot}

\chapter{海の雨 (Sea Rain)}\label{ux6d77ux306eux96e8-sea-rain}

1830年から手つかずの町並みが残る美々津町（みみつまち）。ここは2600年前、日本の初代天皇である神武天皇が日本統一に向けて船出した場所でもある。僕がここに来たのは、単に日本で一番有名なサーフスポットに近かったからだ。でも幸運なことに、そこで「ケイゾウさん」という男性に出会うことができた。

彼は定年退職した中年のサーファーで、世界中からやってくる多くのサーファーたちをガイドしてきた。彼の英語は少し錆びついていて、Google翻訳を通しても理解するのが難しい時があったけれど、僕がこれまで出会った日本人の中で間違いなく最高の人物だった。日焼け対策に余念がない日本では珍しく、彼は太陽の光で真っ黒に日焼けしていた。

それに、長年多くの外国人と接してきたからか、彼の日本的な礼儀正しさは、西洋のサーファー文化にいい意味で「浸食」されているのがわかった。3年前、嵐の中で6メートルの波に飲まれて両膝を脱臼して以来、ずっとサーフィンができなかったと彼は教えてくれた。でもこの週末、僕と一緒に3年ぶりにサーフィンに復帰するというのだ。

とても光栄に思ったけれど、同時に4メートル級の波に連れて行かれるんじゃないかとビクビクもしていた。とにかく、それは宮崎の海で波に顔を激しく打たれた僕の初日だった。控えめに言っても、海は決して穏やかではなかった。

僕は波が自分とボードに襲いかかってくるのを待っていた。ただただ、その波の大きさに圧倒されていた。周りには50人以上のサーファーが波に乗っていて、それはなかなかの光景だった。

30分ほど前に太陽が隠れたかと思うと、分厚い雲が押し寄せてきた。雲は僕たちから方向感覚を奪っていった。1分おきにメガワット級の稲妻が光っていた。そして間もなく、僕らは海に降る朝の雨を迎えた。なんて最高な気分だろう。海の雨の中をパドリングしながら、顔に跳ね返る水滴を感じる。

ありがとう、ケイゾウさん！

\begin{center}\rule{0.5\linewidth}{0.5pt}\end{center}

\bookmarksetup{startatroot}

\chapter{Best meal of my life.}\label{best-meal-of-my-life.}

It was my first and last night, staying at Mimitsu-machi, the old
samurai town of Edo.

The temperature was way too low to surf, I decided to invite Keizo for a
last meal at the neighborhood izakaya. The name was Hinataya.

From outside, it did look like old buildings in the old town, a
19th-century traditional samurai house. But the inside was totally
different.

It was pretty hipster-y, with an old-school wooden pendant clock, a cute
illustration, and 70s rock ballads playing. I wasn't shocked. The cutest
thing were their menus depicting cute little twins and saying something
like, ``Don't mind the two little monsters if you see them running
around.''

I quickly understood that it was not your regular sports izakaya with
drunken dudes sipping on shochu. No, it was a refined and classy
contemporary restaurant run by a young couple that fled from Tokyo's
noise to live there by themselves in a small hotel basement, raising the
two little terrors by the sound of the waves crashing on the shore. I
asked Keizo about his menu and recommendations, and he smiled at me and
said, ``Everything is good'' with a big smile.

So I started decrypting the Japanese menu, Keizo went ahead and ordered
a platter of sashimi and tempura fried fish, to start off. I ordered my
usual drink here in the Kyushu peninsula, and a yuzu cider, a sort of
Japanese soda with a hint of yuzu.

Amazing. My fish platter was already almost finished, and I wasn't
really excited on what to order, but something caught my eye. A whole
``pata de jamón ibérico'' was there laying behind the bar.

I asked the owner, ``Is this from Japan?'' And he said, ``No, no, no,
this is from Spain.'' I started telling him that I was from Spain, but
they were quite confused, since Keizo told them that I was from
Switzerland like 10 minutes ago. Anyway, for a second I thought I was
ordering a platter of ``Berota''. That's how Japanese call ``Bellota'',
because they don't have the L sound, which I found specially hilarious.
But when I saw the price, I quickly backed off. And yeah, I didn't come
all the way here to eat jamón.

What did I order then, you ask me? Well, first I ordered scallops with
asparagus, but of course they were all out of it. That was maybe one of
the best misfortunes I ever had, because after that I ordered the yam
teppanyaki, on the side of their special charcoal-grilled chicken fried
rice. Nothing too fancy.

By that time I already had a bajillion fried rice dishes during my
Japanese adventures, and a billion teppanyaki over my time in the Kansai
area. But oh god, they here it was different. I finally understood the
term ``soul food''. Fried rice is maybe one of the most popular foods
all over Eastern Asia, and also one of the simplest dishes ever. But
man, the charcoal-grilled chicken had so much flavor from the smokiness
and the umaminess. Everything was just so balanced.

And the rice, it didn't look like anything. It was cooked to perfection.
Not too dry, not too oily, not too chunky, not too boiled. The veggies
were barely cooked, just the perfect crunchiness. There were some fresh
radishes cut in very small cubes, and that brought some very much
appreciated freshness. Per-fe-ction. It felt like I was a ten-year-old
boy coming out of school in rural Japan and getting served a bowl of
love by my Japanese grandma. But wait, that was not all.

I was served with this yam teppanyaki in its iron pan. It was smoking
hot. It looked like a Spanish tortilla made love to an American pancake,
but inside it was a pool of cheese, which was not very common to see in
Japan. And it was so fluffy, like if someone gently touched the inside
of your mouth with a sponge. It was very simple, but the Japanese yam,
which is not very sweet like we are used to in Western countries, and
the daikon, Japanese white radishes, cooked to perfection, and the
flavors, again, balanced with the utmost delicacy and respect, just
melted my heart. I kept on making this ``mmmhh'' sound you usually make
to be polite, or when you don't know the language, and you want to
express appreciation, but all of them, every single one, were true, and
couldn't be helped. After savoring each bite of my meal, I smiled and
felt warmness and tenderness inside of me. I couldn't express what I was
feeling in Japanese. It was way more complicated than my basic Japanese.
It's even still hard to describe it in English right now..

\begin{center}\rule{0.5\linewidth}{0.5pt}\end{center}

\bookmarksetup{startatroot}

\chapter{人生最高の食事 (Best meal of my
life)}\label{ux4ebaux751fux6700ux9ad8ux306eux98dfux4e8b-best-meal-of-my-life}

それは、江戸時代の古い武家町である美々津町での、最初で最後の夜だった。

サーフィンをするには気温が低すぎたため、僕はケイゾウさんを近所の居酒屋での最後の晩餐に誘うことにした。お店の名前は「ひなた屋」。

外観は古い町並みに溶け込む、19世紀の伝統的な武家屋敷そのものだった。しかし、一歩中に入ると全くの別世界が広がっていた。

レトロな木製の振り子時計に可愛らしいイラストが飾られ、70年代のロックバラードが流れる、かなりヒップスターな空間だった。でも、驚きはしなかった。一番可愛かったのは、メニューに小さな双子のイラストが描かれていて、「2匹の小さなモンスターが走り回っていても気にしないでくださいね」といったメッセージが添えられていたことだ。

酔っ払ったおじさんたちが焼酎をすすっているような、普通のスポーツ居酒屋ではないことはすぐにわかった。そこは、東京の喧騒から逃れてきた若い夫婦が営む、洗練された上品なモダン・レストランだったのだ。彼らは小さなホテルの地下で、岸に打ち寄せる波の音を聞きながら2人の小さな怪獣を育てている。ケイゾウさんにメニューのおすすめを聞くと、彼は満面の笑みで「全部美味しいよ」と答えた。

僕が日本語のメニューを解読し始めている間に、ケイゾウさんは手始めにと、刺身の盛り合わせと魚の天ぷらを注文してくれた。僕はここ九州でのいつもの飲み物と、ほんのり柚子の風味が効いた日本のソーダ「ゆずサイダー」を頼んだ。

最高だった。魚の盛り合わせがもうすぐなくなりそうで、次に何を頼もうか特にピンときていなかった時、あるものが目に飛び込んできた。カウンターの奥に「イベリコ豚の生ハム」の原木が丸ごと置かれていたのだ。

オーナーに「これ、日本のものですか？」と聞くと、「いやいや、スペイン産ですよ」と返ってきた。僕は自分がスペイン出身だと伝えようとしたけれど、彼らはかなり困惑していた。というのも、つい10分ほど前にケイゾウさんが「彼はスイスから来た」と伝えていたからだ。まあそれはともかく、一瞬「ベロータ」の盛り合わせを頼もうかと考えた。日本語にはLの発音がないから「Bellota」を「ベロータ」と呼ぶんだけど、それが個人的にツボだったのだ。でも、値段を見た瞬間、すぐに身を引いた。それに、わざわざ日本まで来て生ハムを食べる必要もないしね。

じゃあ何を頼んだのかって？
まずはホタテとアスパラガスを注文したんだけど、案の定売り切れだった。でも、それは僕の人生で最高の不運だったかもしれない。なぜなら、その代わりに山芋の鉄板焼きと、特製の炭火焼き地鶏チャーハンを頼むことになったからだ。特に気取った料理ではない。

日本での冒険の中で、僕はすでに星の数ほどのチャーハンを食べてきたし、関西にいた頃には数え切れないほどの鉄板焼きを食べてきた。でも、ああ神様、ここのは全く違った。「ソウルフード」という言葉の意味を、ついに理解した瞬間だった。

チャーハンはおそらく東アジア全域で最も人気のある食べ物の一つであり、同時に最もシンプルな料理でもある。でも、なんということだ。炭火焼きの地鶏には、スモーキーさと旨味がたっぷり詰まっていた。すべてが完璧なバランスだった。

そしてお米。見た目は普通なのに、火の通り具合が完璧だった。パサパサすぎず、油っぽすぎず、ダマにもなっておらず、柔らかすぎない。野菜はほんの少し火が通っただけで、完璧なシャキシャキ感。そこへ小さなサイコロ状に切られた新鮮な大根が入っていて、それがまたありがたい清涼感をもたらしていた。完・璧、だ。まるで自分が日本の田舎の学校から帰ってきた10歳の少年で、日本のおばあちゃんから愛情たっぷりのご飯を出してもらったような気分だった。でも待ってほしい、それだけじゃない。

鉄板に乗ったまま運ばれてきた山芋の鉄板焼き。湯気を立てて熱々だった。見た目はスペイン風オムレツとアメリカンパンケーキが恋に落ちて生まれたような感じだったけれど、中には日本では珍しいチーズの海が広がっていた。そして、まるで誰かがスポンジで口の中を優しく触っているかのようにフワフワだった。とてもシンプルだけれど、欧米の芋のように甘すぎない日本の山芋と、完璧に火が通った大根、そして最高の繊細さと敬意をもってバランスが取られたその風味が、僕の心を溶かした。礼儀として、あるいは言葉がわからない時に感謝を伝えるために出す「んーー」というあの唸り声を何度も漏らしてしまったけれど、そのすべてが、一つ残らず心からのものであり、どうしようもないほど溢れ出てしまったものだった。

ひと口ひと口を味わった後、僕は微笑み、自分の中に温かさと優しさを感じていた。自分が感じていることを日本語で表現することはできなかった。僕の拙い日本語レベルを遥かに超えるほど、それは複雑な感情だった。今、英語で書き表すことすら難しいくらいに。

\bookmarksetup{startatroot}

\chapter{The Weight of Leaving}\label{the-weight-of-leaving}

Naha feels different when you know you're about to leave. Louder.
Closer. Harder to hold at arm's length.

The city hums constantly, traffic, voices, engines, the low mechanical
breath of the nearby army base. Jets cut the sky without warning,
shaking windows, rattling thoughts. The sound never quite fades, it just
shifts position, like something pacing in the background.

Everything is tinted in soft contradiction. Pastel green cars parked
along narrow streets. Houses painted in faded blues and mint tones,
sun-worn. And then the neon, sharp and restless, slicing through the
night, convenience stores glowing like islands of artificial daylight.
The immensity of the Don Quijote stores takes over the night, their
endless aisles and looping music bleeding into the streets.

I walk without direction, overwhelmed by it all. The density. The
lights. The noise. The feeling that the city is both holding me and
pushing me away. Naha doesn't ask you to stay still. It keeps moving,
insists on it, even when you don't have the strength to follow.

Leaving Okinawa from Naha feels like being gently but firmly pulled back
into the world. And as much as the city overwhelms me, it's that excess,
that noise, that uneasy mix of softness and force, that makes saying
goodbye hurt more than I expected..

\begin{center}\rule{0.5\linewidth}{0.5pt}\end{center}

\bookmarksetup{startatroot}

\chapter{手放すには近すぎる}\label{ux624bux653eux3059ux306bux306fux8fd1ux3059ux304eux308b}

那覇は、これから離れると分かっているとき、少し違って見える。より騒がしく、より近く、距離を保つのが難しくなる。

この街は絶えずうなり続けている。車の音、人の声、エンジンの気配、近くの軍事基地から届く低い機械音の呼吸。予告もなくジェット機が空を切り裂き、窓を揺らし、思考を乱す。その音は完全に消えることがない。ただ位置を変えるだけで、背後を歩き続ける何かのように残り続ける。

すべてが、やわらかな矛盾の中に染まっている。細い路地に並ぶ淡い緑の車。色あせた青やミント色の家々。そこに、夜を切り裂くようなネオンの鋭さが重なる。コンビニは人工的な昼の光のように光り続け、ドン・キホーテの巨大な存在感が夜を覆い尽くす。終わりのない棚と、繰り返される音楽が、街へと滲み出していく。

行き先もなく歩く。ただ圧倒されている。密度に。光に。音に。この街が自分を抱え込みながら、同時に押し返してくるような感覚に。那覇は静止を許さない。たとえそれに従う力が残っていなくても、動き続けることを要求してくる。

那覇から沖縄を離れるというのは、世界へとやわらかく、しかし確実に引き戻されていく感覚に近い。そして、この街の過剰さに押し流されながらも、その過剰さ、その音、その不安定なやさしさと暴力の混ざり合いこそが、想像以上に別れを痛くしている。

\bookmarksetup{startatroot}

\chapter{Tokashiku beach}\label{tokashiku-beach}

C'est un canal d'eau morte entre les berges de la rivière et l'océan
Pacifique, ou des milliers de bernard-l'hermite dressés devant leur trou
battaient l'air de leur pince droite dans des gestes racoleurs.

\begin{center}\rule{0.5\linewidth}{0.5pt}\end{center}

\subsection{渡嘉敷ビーチ}\label{ux6e21ux5609ux6577ux30d3ux30fcux30c1}

それは、川岸と太平洋のあいだに挟まれた死んだ水のような水路だった。そこでは何千ものヤドカリが自分の穴の前に立ち、右のはさみを空中に向かって振り上げていた。まるで何かを呼び込むような、軽薄な仕草で。

\begin{center}\rule{0.5\linewidth}{0.5pt}\end{center}

\subsection{Normal}\label{normal}

Normal parfois je repense à mon adolescence et j'en dors mal. Je pensais
que c'était normal. Il y a des choses dans ma tête que je dirais jamais
à mes parents. C'est pas en écrivant quoi que ce soit que ça risque de
changer. J'effacerai pas mes problèmes comme ça je le sais depuis des
années. J'ai les mêmes habitudes depuis que je suis enfant. Je suis doué
pour fuir la réalité et pour faire semblant. Ne pas voir le cœur des
adultes qui saignent. C'est ça de grandir sans avoir des parents qui
s'aiment. Pas un câlin pas un bisou je ne réalisais pas. J'ai découvert
l'amour à la télé, au cinéma. Pas étonnant que je confonde mes romans
avec la vie. Que je n'arrive pas à dire'' je t'aime'' à une
fille\ldots{}

\subsection{ノーマル}\label{ux30ceux30fcux30deux30eb}

ときどき、自分の思春期のことを思い出して眠れなくなることがある。それが普通だと思っていた。頭の中には、両親に絶対に言わないことがいくつもある。何かを書いたところで、それで変わるとも思えない。問題が消えることはないと、ずっと前から分かっている。

子どもの頃から同じ癖を持っている。現実から逃げるのが得意で、何もないふりをするのが上手い。大人たちの心が血を流していることにも気づかないふりをする。

愛し合わない両親のもとで育つというのはそういうことだ。抱きしめも、キスもないまま、それが普通だと思っていた。愛というものはテレビや映画で知った。だからこそ、小説と現実を混同してしまう。女の子に「愛してる」と言えないのも、そのせいかもしれない。

\begin{center}\rule{0.5\linewidth}{0.5pt}\end{center}

\subsection{Scorpion Fish (By Nicolas Bouvier translated from
French)}\label{scorpion-fish-by-nicolas-bouvier-translated-from-french}

\subsection{スコーピオン・フィッシュ
(ニコラ・ブーヴィエ著、フランス語からの翻訳)}\label{ux30b9ux30b3ux30fcux30d4ux30aaux30f3ux30d5ux30a3ux30c3ux30b7ux30e5-ux30cbux30b3ux30e9ux30d6ux30fcux30f4ux30a3ux30a8ux8457ux30d5ux30e9ux30f3ux30b9ux8a9eux304bux3089ux306eux7ffbux8a33}

The more you bring to the island, the more it gives back to you.
島に持ち込めば持ち込むほど、島はそれ以上のものを返してくれる。

He who does not accept to begin by learning to live with less here is
certain to waste his time.
ここで、より少ないもので生きることを最初に学ぼうとしない者は、間違いなく時間を無駄にすることになる。

Traveling: putting your head on the chopping block a hundred times, and
a hundred times retrieving it from the basket to find it almost
unchanged. We still hoped for a miracle, even though we shouldn't expect
any other than this wearing down and this erosion of life with which we
have an appointment, and before which we so wrongly rear up. I travel to
learn, and no one had taught me what I discovered here.
旅とは、100回首切り台に頭を乗せ、100回その首を籠から拾い上げ、それでもほとんど変わっていない自分の顔を見つけることだ。私たちはそれでも奇跡を期待してしまうが、私たちが直面することになる生命の摩耗や浸食、私たちがひどく間違って抵抗してしまうその現象以上に、別の奇跡を待つべきではないのだ。私は学ぶために旅をしているが、私がここで発見したことを教えてくれた者は誰もいなかった。

\subsection{Puesta de Sol (spanish)}\label{puesta-de-sol-spanish}

Entre los árboles pude verlas nuevas luces del amanecer punto entre las
branchas podíamos ver el sol naranja del atardecer.

\begin{center}\rule{0.5\linewidth}{0.5pt}\end{center}

\subsection{日没}\label{ux65e5ux6ca1}

木々の間から、新しい朝の光が見えた。枝のあいだに、オレンジ色の夕日の太陽が覗いていた。

\subsection{forest Bridges.}\label{forest-bridges.}

From the ferry you can see those three huge Bridges seemingly coming and
going from nowhere inside the bushy, untouched Forest Hills of
tokashiki. After a bit of determination and a lot of sweat you can get
to the top of nishiyama between parentheses West mountain to observe one
of the most wonderful sites I have ever seen. A red hot perfect Circle
falling behind the slime is Shadows of the karama archipelago

--

\subsection{森の橋}\label{ux68eeux306eux6a4b}

フェリーからは、三本の巨大な橋が見える。それは、まるで何もない森の中から現れては消えていくように見えた。決意と汗を重ねながら西山（西の山）へ登ると、これまで見た中でも最も美しい光景に出会うことができる。

赤く燃える完璧な円が、影のように広がる慶良間諸島の向こうへ沈んでいく。

\subsection{Ravens}\label{ravens}

Even if ravens are magnificent, they are doomed to die alone, whereas
pigeons can share everything.

\begin{center}\rule{0.5\linewidth}{0.5pt}\end{center}

\subsection{カラス}\label{ux30abux30e9ux30b9}

ワタリガラスはたとえ壮麗であっても孤独に死ぬ運命にあるが、ハトはすべてを分かち合うことができる。

15th of November, 2025.

\bookmarksetup{startatroot}

\chapter{Ferry Thoughs}\label{ferry-thoughs}

After seeing the sun slip behind a the Sakurajima volcano, from the
Kagoshima ferry's upper deck, feels like a quiet moment, stretching
through eternity. The sea is calm, the sky slowly closing.

It's beautiful, but it also stings a little. That mix of leaving
something behind and moving towards something unknown, made me think of
migrants seeing their coastline fade for the last time, maybe the same
light on the water, but heavy on the chest. Hope in one hand, fear on
the other.

A silent promise that whatever waits across the sea is worth the risk,
even if they are not sure yet.

Just a few minutes of sunset, but so many stories that have crossed that
horizon.

\begin{figure}[H]

{\centering \pandocbounded{\includegraphics[keepaspectratio]{.attachments/7b80cdf61657afa216a8a87a15e95990814fe9c2.jpeg}}

}

\caption{image}

\end{figure}%

\begin{figure}[H]

{\centering \pandocbounded{\includegraphics[keepaspectratio]{.attachments/ac89d37b4e098be57fca047f4abdf4d358880df3.jpeg}}

}

\caption{image}

\end{figure}%

\begin{center}\rule{0.5\linewidth}{0.5pt}\end{center}

\subsection{フェリーの思考}\label{ux30d5ux30a7ux30eaux30fcux306eux601dux8003}

鹿児島のフェリーの上甲板から、桜島の火山の向こうへ太陽が沈んでいくのを見ていた。その光景は、静かで、永遠の中に引き延ばされていくような瞬間だった。海は穏やかで、空はゆっくりと閉じていく。

美しい光景だったが、少しだけ胸を刺すものがあった。何かを置き去りにする感覚と、未知へ向かう感覚が同時にある。その感覚が、海岸線が最後に消えていくのを見つめる移民たちの姿を思い起こさせた。同じ水面の光かもしれないが、胸の奥には重さがある。片手には希望、もう片手には恐れ。

海の向こうに何が待っているのか分からなくても、それでも進む価値があるという、静かな約束。

ただ数分の夕日。それでも、その地平線を越えてきた無数の物語がある。


\backmatter


\end{document}
